\documentclass[10pt,a4paper]{amsart}
\usepackage[margin=19mm]{geometry}
\usepackage{amsmath,amssymb,mathtools}
\usepackage[scr=rsfs,cal=euler]{mathalfa}
\usepackage{xfrac}
\usepackage[unicode,hidelinks,bookmarks=false]{hyperref}
\newcommand{\ie}{i.e.\ }
\newcommand{\cf}{cf.\ }
\newcommand{\resp}{respectively\ }
\newcommand{\Subsection}{\subsection}
\providecommand{\nobreakdash}{\nobreak\hbox{-}}
\setlength{\emergencystretch}{2em}
\multlinegap0pt
\usepackage{bm}
\usepackage{bbm}
\usepackage{dsfont}
\usepackage{xspace}
\usepackage{cancel}
\usepackage{mathtools}
\usepackage{amsthm}
\makeatletter
\@ifundefined{theorem}{\newtheorem{theorem}{Theorem}}{}
\@ifundefined{lemma}{\newtheorem{lemma}[theorem]{Lemma}}{}
\makeatother
\makeatletter
\@ifundefined{proofthm}{\newenvironment{proofthm}[1]{%
  \begin{proof}[Proof of Theorem~#1]%
}{%
  \end{proof}%
}}{}
\makeatother
\title[A Cheng-type Eigenvalue-Comparison Theorem for the Hodge Lap]{A Cheng-type Eigenvalue-Comparison Theorem for the Hodge Laplacian}
\author{Anusha Bhattacharya, Soma Maity}
\date{Source: arXiv:2603.10633v1 (CC BY 4.0)}
\begin{document}
\maketitle

\noindent\emph{Source and licence.} A Cheng-type Eigenvalue-Comparison Theorem for the Hodge Laplacian. Version arXiv:2603.10633v1 (CC BY 4.0), distributed under the Creative Commons Attribution 4.0 International licence, as recorded in the arXiv licence metadata for this exact version and shown on the abstract page https://arxiv.org/abs/2603.10633v1. Full rights notice: licenses/arxiv-2603.10633-CC-BY-4.0.txt.\par\smallskip

\noindent\textit{Editorial scope.} Counted unit: the Theorem whose source label is \texttt{thm:volume} in the article (source0.tex lines 306--316, source label \texttt{thm:volume}), delivered together with the article's own complete original proof of it (source0.tex lines 317--380, one \texttt{proof} environment of 64 lines). No mathematics is written, repaired or re-derived: statement and proof are the article's own text line for line. Required context retained uncounted: lem:upper bound dirichlet eigenvalue (lines 201--205). Every definition, display and label of those lines is retained unchanged. Omitted from this package and not redistributed: 1--200, 206--305, 381--483. Nothing inside a retained excerpt is omitted or altered: every retained body is delivered exactly as the article's lines are written, so no omission receipt of any kind accompanies this package. Citations: the retained excerpts carry no citation command at all, so no bibliography entry of the article is transcribed and no reference is redistributed. Reference adaptation: none was needed and none was made; every reference inside the delivered unit resolves inside it, so no unresolved reference or citation is produced. Printed slips kept literal: no printed word, symbol, hypothesis or number of the counted statement or of its proof was altered. Layout and class: the article's own class is \texttt{amsart} with packages including mathtools, xcolor, geometry, latexsym, mathrsfs, enumerate, fullpage, graphicx, amsmath, amsfonts, amssymb, amsthm, tikz-cd, biblatex; the wrapper is the lane's pinned \texttt{amsart} editorial wrapper at 10pt on A4, which restates the theorem-like environments the retained text needs and adds the article's own macro definitions that the retained text needs (none), with extra wrapper packages (bm, bbm, dsfont, xspace, cancel, mathtools). The delivered numbering is the unit's own. Assets: no retained span contains a figure environment, a graphics-inclusion command, a PDF-page inclusion, a table or a TikZ picture; no image file is copied into this package, the package declares no asset dependency and no image file is redistributed with it. Licence: version arXiv:2603.10633v1 of this article is distributed under the Creative Commons Attribution 4.0 International licence, as recorded in the arXiv licence metadata for this exact version and shown on the abstract page https://arxiv.org/abs/2603.10633v1. A full rights notice is delivered at licenses/arxiv-2603.10633-CC-BY-4.0.txt. Characters: every retained excerpt is pure ASCII, so the delivered engine, XeLaTeX with its default Latin Modern text font, reports no missing character.
\begin{lemma}    \label{lem:upper bound dirichlet eigenvalue}
Let $(M,g)\in \mathcal{M}(n,\xi,r_0)$ and $x\in M.$ For $0<\varepsilon\le r_H$,  let $B=B(x, \varepsilon)$ and $B_\xi(\varepsilon)$ be a ball of radius $\varepsilon$ in the space of constant curvature $\xi.$ Then   
\[\lambda_{0,p}^D(B)\le 2^{2p+1}\lambda_0^D(B_\xi(\varepsilon))\]
where $\lambda_0^D(B)$ denotes the first Dirichlet eigenvalue of the Laplace-Beltrami operator acting on functions.
\end{lemma}

\begin{theorem}
\label{thm:volume}
 Let $0\le r_H\le \frac{1}{\sqrt{|\xi|}}$ and $(M,g)$ be a closed Riemannian $n$-manifold with volume $V$. Assume that the Ricci curvature satisfies $\mathrm{Ric}_g \ge (n-1)\xi .$ Then   for $n>1$,

\[
\lambda_{k,p}
\le
2^{2p+n+5}\cdot (k+1)\cdot \left(\frac{\alpha_n}{V}\right)^\frac{2}{n},\qquad \mathrm{for }\hspace{1mm} k>\frac{2^nV}{\alpha_{n-1}r_H^n}.
\]
where $\alpha_n=\mathrm{vol}(S^n).$
\end{theorem}

\begin{proof}
We choose \[\varepsilon
=
\frac{2}{\sqrt{|\xi|}}\operatorname{arcsinh}\left[\left(\frac{ V|\xi|^\frac{n}{2}}{\alpha_n(k+1)}\right)^{1/n}\right].
\]  

From Bishop's theorem, 
\begin{equation}
\label{eq:volume_ratio}
  |X(\varepsilon)|\ge \frac{V}{\sup_{x\in X(\varepsilon)}\{B(x,\varepsilon/2)\}}\ge \frac{V}{V_{\xi}(\varepsilon/2)}.  
\end{equation}
\begin{align}
\mathrm{Vol}_{\xi}(\varepsilon/2)
&=
\alpha_{n-1}
\int_{0}^{\varepsilon/2}
\left(
\frac{\sinh(\sqrt{|\xi|}t)}{\sqrt{|\xi|}}
\right)^{n-1}
dt \nonumber\\
&=
\frac{\alpha_{n-1}}{|\xi|^{n/2}}
\int_0^{\varepsilon \sqrt{|\xi|}/2}
\sinh^{n-1}(t)\,dt \nonumber\\
&\le
\frac{\alpha_{n-1}}{|\xi|^{n/2}}
\left(
\sinh\left(\frac{\varepsilon \sqrt{|\xi|}}{2}\right)
\right)^n \nonumber
\end{align}

Substituting the above bound and chosen $\varepsilon$ in (\ref{eq:volume_ratio}),

\begin{align}
|X(\varepsilon)|\ge 
\frac{V |\xi|^{n/2}}
{\alpha_{n-1}
\left(
\sinh\left(\frac{\varepsilon \sqrt{|\xi|}}{2}\right)
\right)^n} \ge k+1.
\end{align}
Now $y\le \sinh{y}\implies \frac{1}{y^n}\ge \frac{1}{(\sinh y)^n}$. Hence for $k>\frac{2^nV}{\alpha_{n-1}r_H^n}$, we have $k>\frac{V|\xi|^\frac{n}{2}}{\alpha_{n-1}(\sinh(r_H\sqrt{|\xi|}/2))^n}.$ Hence,  $\varepsilon<r_H$ for $k>\frac{2^nV}{\alpha_{n-1}r_H^n}$ . Using Lemma \ref{lem:upper bound dirichlet eigenvalue}, for $f:B(x,\varepsilon)\rightarrow\mathbb{R}$ defined as $f(y)=1-\frac{1}{\varepsilon}d(x,y)$. Then $|df|\ge \frac{1}{\varepsilon}$ and $|f|\ge\frac{1}{2}$ in $B\left(x,\frac{\varepsilon}{2}\right)$. Hence, using $\omega=f\omega_0$ as the test form in the lemma and using Bishop-Gromov volume ratio, we have
\begin{equation}
    \label{eq:lambda upper bound for volume}
    \lambda_{k,p}\le 2^{2p+3}\frac{V_\xi(\varepsilon)}{\varepsilon^2 V_\xi(\frac{\varepsilon}{2})}, \qquad \mathrm{for}\hspace{1mm} k>\frac{2^nV}{\alpha_{n-1}r_H^n}.
\end{equation}
Let $x^n=\frac{ V|\xi|^\frac{n}{2}}{\alpha_n(k+1)}$. Then \begin{align}
    \frac{V_\xi(\varepsilon)}{ V_\xi(\frac{\varepsilon}{2})} &\le \left[\frac{\sinh\left(\varepsilon\sqrt{|\xi|}\right)}{\sinh\left(\varepsilon\sqrt{|\xi|}/2\right)}\right]^n \nonumber\\
    &\le 2^n\left[\cosh\left(\varepsilon\sqrt{|\xi|}/2\right)
    \right]^n \nonumber \\
    &\le 2^n\left[\cosh\left(\sinh^{-1}x\right)\right]^n \nonumber\\
    &\le 2^n\left(1+x\right)^n.
\end{align}
Since $\sinh^{-1}x\ge \frac{x}{1+x}$, substituting $\varepsilon$ in (\ref{eq:lambda upper bound for volume})  gives
\begin{equation}
\label{eq:upper bound after sinh substitution}
    \lambda_{k,p}\le 2^{2p+3}|\xi|\frac{(1+x)^{n+2}}{x^2}.
\end{equation}
Now $r_H\le \frac{1}{\sqrt{|\xi|}}$implies $\sinh \left(r_H\sqrt{|\xi|}/2\right)<1$. Hence,  $k>\frac{2^nV}{\alpha_{n-1}r_H^n}>\frac{V|\xi|^\frac{n}{2}}{\alpha_{n-1}(\sinh(r_H\sqrt{|\xi|}/2))^n}> \frac{V|\xi|^\frac{n}{2}}{\alpha_{n-1}}$ which implies $x<1.$\\

Substituting in (\ref{eq:upper bound after sinh substitution}),
\[ \lambda_{k,p}\le 2^{2p+n+5}\cdot (k+1)\cdot \left(\frac{\alpha_n}{V}\right)^\frac{2}{n}.\]

\end{proof}

\end{document}
